\documentclass[10pt,a4paper]{amsart}
\usepackage{amsmath,amssymb,amsthm,graphicx}
\usepackage{verbatim}
\newenvironment{reference}{\comment}{\endcomment}
%\newenvironment{reference}{}{}
\newenvironment{slogan}{\comment}{\endcomment}
\newenvironment{history}{\comment}{\endcomment}
% For commutative diagrams we use Xy-pic
\usepackage[all]{xy}

% We use 2cell for 2-commutative diagrams.
\xyoption{2cell}
% Theorem environments.
%
\theoremstyle{plain}
\newtheorem{theorem}[subsection]{Theorem}
\newtheorem{proposition}[subsection]{Proposition}
\newtheorem{lemma}[subsection]{Lemma}

\theoremstyle{definition}
\newtheorem{definition}[subsection]{Definition}
\newtheorem{example}[subsection]{Example}
\newtheorem{exercise}[subsection]{Exercise}
\newtheorem{situation}[subsection]{Situation}

\theoremstyle{remark}
\newtheorem{remark}[subsection]{Remark}
\newtheorem{remarks}[subsection]{Remarks}

\numberwithin{equation}{subsection}

% Macros
%
\def\lim{\mathop{\mathrm{lim}}\nolimits}
\def\colim{\mathop{\mathrm{colim}}\nolimits}
\def\Spec{\mathop{\mathrm{Spec}}}
\def\Hom{\mathop{\mathrm{Hom}}\nolimits}
\def\Ext{\mathop{\mathrm{Ext}}\nolimits}
\def\SheafHom{\mathop{\mathcal{H}\!\mathit{om}}\nolimits}
\def\SheafExt{\mathop{\mathcal{E}\!\mathit{xt}}\nolimits}
\def\Sch{\mathit{Sch}}
\def\Mor{\mathop{\mathrm{Mor}}\nolimits}
\def\Ob{\mathop{\mathrm{Ob}}\nolimits}
\def\Sh{\mathop{\mathit{Sh}}\nolimits}
\def\NL{\mathop{N\!L}\nolimits}
\def\CH{\mathop{\mathrm{CH}}\nolimits}
\def\proetale{{pro\text{-}\acute{e}tale}}
\def\etale{{\acute{e}tale}}
\def\QCoh{\mathit{QCoh}}
\def\Ker{\mathop{\mathrm{Ker}}}
\def\Im{\mathop{\mathrm{Im}}}
\def\Coker{\mathop{\mathrm{Coker}}}
\def\Coim{\mathop{\mathrm{Coim}}}

% Boxtimes
%
\DeclareMathSymbol{\boxtimes}{\mathbin}{AMSa}{"02}

%
% Macros for moduli stacks/spaces
%
\def\QCohstack{\mathcal{QC}\!\mathit{oh}}
\def\Cohstack{\mathcal{C}\!\mathit{oh}}
\def\Spacesstack{\mathcal{S}\!\mathit{paces}}
\def\Quotfunctor{\mathrm{Quot}}
\def\Hilbfunctor{\mathrm{Hilb}}
\def\Curvesstack{\mathcal{C}\!\mathit{urves}}
\def\Polarizedstack{\mathcal{P}\!\mathit{olarized}}
\def\Complexesstack{\mathcal{C}\!\mathit{omplexes}}
% \Pic is the operator that assigns to X its picard group, usage \Pic(X)
% \Picardstack_{X/B} denotes the Picard stack of X over B
% \Picardfunctor_{X/B} denotes the Picard functor of X over B
\def\Pic{\mathop{\mathrm{Pic}}\nolimits}
\def\Picardstack{\mathcal{P}\!\mathit{ic}}
\def\Picardfunctor{\mathrm{Pic}}
\def\Deformationcategory{\mathcal{D}\!\mathit{ef}}
\title{The Euler sequence on projective space, with the two maps written out on the standard opens}
\author{The Stacks Project Authors}
\date{stacks-project:derham.tex:local-snapshot-20261012; GNU FDL 1.2 or later}
\begin{document}
\maketitle
\noindent\textit{Editorial note.} All delivered mathematical text is the author's own, copied from the pinned Stacks Project chapter file derham.tex; no statement or proof step was rewritten, reordered or removed.\par\medskip
\noindent\textit{Editorial note.} Preamble provenance: the pinned chapter file opens with \textbackslash{}input\{preamble\} and carries no preamble text of its own, so the project's single Stacks-wide preamble staged at sources/stacks-derham/preamble.tex is excerpted verbatim as three spans, lines 8-12 (verbatim package and the project's own reference/slogan/history display environments), lines 14-18 (Xy-pic setup) and lines 154-217 (theorem declarations and mathematical macros). The wrapper class and the amsmath/amssymb/amsthm/graphicx base packages are supplied by the assembler. No macro was invented or redefined.\par\medskip
\noindent\textit{Editorial note.} Same-chapter pointers are located at their \textbackslash{}label lines inside the pinned chapter file; cross-chapter pointers are located in the project's own tag list staged in the snapshot (sources/stacks-chow/tags.txt), at the verbatim lines carrying their canonical tags.\par\medskip
\noindent\textit{Allowed standard prerequisites (editorial).}\begin{itemize}
\item Projective space as Proj of a polynomial ring with its standard opens. Same chapter derham.tex, the discussion preceding the lemma (used at proof lines 1602-1617). It fixes \$\textbackslash{}mathbf\{P\}\textasciicircum{}n\_A = \textbackslash{}operatorname\{Proj\}(A[T\_0, \textbackslash{}ldots, T\_n])\$ and the standard opens \$D\_+(T\_i) = \textbackslash{}Spec(A[T\_0/T\_i, \textbackslash{}ldots, T\_n/T\_i])\$ on which the maps are written.
\item The tautological quotient and the identifications of \$\textbackslash{}mathcal\{O\}(-1)\$ and \$\textbackslash{}mathcal\{O\}\$ on the standard opens. Same chapter derham.tex, the description of the relevant invertible sheaves used at proof lines 1606-1608 and 1624. They let \$\textbackslash{}mathcal\{O\}(-1)\textasciicircum{}\{\textbackslash{}oplus n+1\}\$ be written as \$\textbackslash{}bigoplus\_j \textbackslash{}mathcal\{O\}(-1)\textbackslash{}mathrm\{d\}T\_j\$ and make sense of the coefficients \$1/T\_i\$ and \$T\_j\$ appearing in the two maps.
\end{itemize}\medskip
\begin{lemma}
\label{lemma-euler-sequence}
There exists a short exact sequence
$$
0 \to \Omega \to \mathcal{O}(-1)^{\oplus n + 1} \to \mathcal{O} \to 0
$$
\end{lemma}

\begin{proof}
To explain this, we recall that
$\mathbf{P}^n_A = \text{Proj}(A[T_0, \ldots, T_n])$,
and we write symbolically
$$
\mathcal{O}(-1)^{\oplus n + 1} =
\bigoplus\nolimits_{j = 0, \ldots, n} \mathcal{O}(-1) \text{d}T_j
$$
The first arrow
$$
\Omega \to
\bigoplus\nolimits_{j = 0, \ldots, n} \mathcal{O}(-1) \text{d}T_j
$$
in the short exact sequence above
is given on each of the standard opens
$D_+(T_i) = \Spec(A[T_0/T_i, \ldots, T_n/T_i])$
mentioned above by the rule
$$
\sum\nolimits_{j \not = i} g_j \text{d}(T_j/T_i)
\longmapsto
\sum\nolimits_{j \not = i} g_j/T_i \text{d}T_j
- (\sum\nolimits_{j \not = i} g_jT_j/T_i^2) \text{d}T_i
$$
This makes sense because $1/T_i$ is a section of $\mathcal{O}(-1)$
over $D_+(T_i)$. The map
$$
\bigoplus\nolimits_{j = 0, \ldots, n} \mathcal{O}(-1) \text{d}T_j
\to
\mathcal{O}
$$
is given by sending $\text{d}T_j$ to $T_j$, more precisely, on
$D_+(T_i)$ we send the section $\sum g_j \text{d}T_j$ to
$\sum T_jg_j$. We omit the verification that this produces
a short exact sequence.
\end{proof}

\end{document}
